% Fresh reconstruction from retained formulas and independently rederived algebra.
% This is NOT a recovered byte stream of r1_unbounded_q_full_conductor.tex.
\section{Reconstructed full R1 conductor in the unbounded source-ratio sector}
\label{sec:fn-reconstructed}

\paragraph{Provenance and precise dependency.}
This is a newly written reconstruction of the late R21 conductor, not a
claim of byte identity with the lost final file.  The complete parent,
the radial Gram matrices, and the normal formulas were retained in the
conversation; every displayed polynomial identity below is checked
again against the recovered original source equations.  The companion
\texttt{check\_fn\_reconstruction.py} is also new.  The source dictionary
is the independently rederived
\texttt{c144/rederived/r1\_qo\_compact\_source.py}.
The uniform source bounds and the scalar allocation stated in the next
paragraph are the explicit interface to the reconstructed actual-source
scalar lemma (RS).  The source input is now supplied by the freshly reconstructed
\texttt{c144/rederived/r1\_qo\_actual\_source\_and\_scalar.tex},
RSR1--RSR29, with its independent current review.  The full-normal
argument below has now received an independent current analytical review
by the lead agent, separate from the review of its scalar input.  Neither
the old admission record nor the old checksum substitutes for it.

\paragraph{Fields, source scope, and notation.}
Let $K_c,K_t\supset\mathbb R$ be real closed fields, with bounded rings
$B_c,B_t$ and infinitesimal ideals $\mathfrak m_c,\mathfrak m_t$.
The ring of certificate expressions is
$\mathcal T=K_c\otimes_{\mathbb R}K_t$, and its bounded coefficient ring
is $\mathcal B=B_c\otimes_{\mathbb R}B_t$.  A bounded coefficient with
zero real residue belongs to
$\mathfrak n=\mathfrak m_c\otimes B_t+B_c\otimes\mathfrak m_t$.
Every radical or inverse used below is pure source, pure testing, or
an explicitly separated product of such factors.  No order on
$\mathcal T$ and no inverse of a mixed expression is used.

Use the positive compact source coordinates
\[
 x=jB,\qquad y=TB,\qquad q_z=jT,\qquad
 0<p<\min\{1,j\},\qquad 0<H<1,                         \tag{RFN1}
\]
where $B,T,x,y,q_z\in\mathfrak m_c$.  These are the compact coordinates
of the original R1 sector with unbounded source ratio.  The original
strict source gaps remain in force and $n\geq0$, including $n=0$.
For reference, their native substitution is
\[
 r=\frac{j-p}{j+px},\quad
 h=\frac{1-HT}{1+HTx},\quad
 a_s=\frac{j+px}{1+HTx},\quad
 B_s=\frac{B(1+HTx)}{1+pB},\quad
 t_s=\frac{T(1+pB)}{1+HTx}.                              \tag{RFN2}
\]
The six weights $n,m,C,A,D,\alpha$ are always the original source
weights at (RFN2), rather than weights of an auxiliary stationary
source.  Put $\mu=m(1-p)$, $h_0=\mu+\tfrac12$,
$k=(j+1-p)/(j+1)$, and $\eta_*=Bj(j+1)/p$.
In this source sector RS proves
\begin{gather}
 \tfrac12<k<1,\quad \eta_*=O(1),\quad
 \mu=O(p),\quad h_0\asymp1,\quad A=O(T),\quad
 \alpha\asymp B^2(j+1),                                  \tag{RFN3a}\\
 TCH\sim k/2,\qquad T(D+\tfrac12)(1-H)\sim k/2,\qquad
 \frac{n(j-p)}j=O(1),                                   \tag{RFN3b}\\
 f_c\asymp Tx^2,\qquad q_c\asymp Tx,\qquad
 r_c\asymp TjB^2(1-p),                                  \tag{RFN3c}\\
 \eta_*\in\mathfrak m_c\quad\Longrightarrow\quad
 n\sim\frac{p}{2(j-p)(j+1)}>0.                           \tag{RFN3d}
\end{gather}
Here $\asymp$ means that the ratio is a positive bounded unit, and
$\sim$ means that the ratio has residue one.  All estimates are uniform
at the independent openings $H$, $1-H$, $j-p$, and $1-p$.
In particular $B(j+1)=B+x$ and $\alpha$ are infinitesimal.
No comparison between source and testing infinitesimals is allowed.

\paragraph{The local conclusion.}
For all positive independent $u,v\in\mathfrak m_t$, normalize the pure
testing coordinate $w=(W+uY)/(1+u^2)$ to one and allow the other pure
testing coordinate $z=(Z-uX)/(1+u^2)>0$ to be arbitrary.  The conclusion
is a finite tensor sum of squares for the complete original parent,
with both state normals unrestricted.  In particular $z$ need not be
bounded.  The separate full-state sign reduction and outer-testing join
are needed to pass from this local conclusion to the whole R1 row.

\subsection{The complete parent and the correlated source mass}

Write $D_u=1+u^2$, $D_v=1+v^2$.  The original quarter-parent is
\begin{align}
 \mathcal P={}&m(vW-X+uvY-uZ)^2+nD_uD_vX^2
  +CD_v(uW-Y)^2+AD_uD_vY^2\\
 &+(D+\tfrac12)(uW+uvX-Y-vZ)^2
  +(\alpha+\tfrac12)D_u(vY-Z)^2\\
 &+\tfrac12D_v(X+uZ)^2-\tfrac12D_uD_vZ^2.               \tag{RFN4}
\end{align}
Use the two mixed normal rows
\[
 N_0=\frac{X+uZ-pD_uv}{D_u},\qquad
 M_0=\frac{uW-Y-HD_uvz}{D_u}.                            \tag{RFN5}
\]
Their exact inverse, with $w=1$, is
\[
 \begin{aligned}
 X&=N_0+pv-uz,&Z&=z+uN_0+puv,\\
 W&=1+u(M_0+Hvz),&Y&=u-M_0-Hvz.
 \end{aligned}                                           \tag{RFN6}
\]
Only the pure testing unit $D_u$ is inverted.  Retain the nonnegative
packet $mpD_u^2N_0^2$ and define
$G=(\mathcal P-mpD_u^2N_0^2)/D_u$.  Direct substitution gives the
\emph{entire} expression
\begin{align}
 G={}&\mu D_u\{N_0^2-2vN_0+(1-p)v^2\}
       +nD_v(N_0+pv-uz)^2\\
 &+CD_uD_v(M_0+Hvz)^2+A D_v(u-M_0-Hvz)^2\\
 &+(D+\tfrac12)D_u(M_0-(1-H)vz)^2\\
 &+(\alpha+\tfrac12)
   \{u(1-p)v-(1+Hv^2)z-uN_0-vM_0\}^2\\
 &+\tfrac12D_uD_v(N_0+pv)^2
       -\tfrac12D_v(z+uN_0+puv)^2.                      \tag{RFN7}
\end{align}
Thus $G$ is exactly quadratic in $(N_0,M_0)$:
\[
 G=F+2L_NN_0+2L_MM_0+H_NN_0^2+H_MM_0^2
                  +2(\alpha+\tfrac12)uvN_0M_0,          \tag{RFN8}
\]
where $F=G|_{N_0=M_0=0}$ and
\[
 \begin{split}
 H_N&=h_0D_u+\tfrac12v^2+\alpha u^2+nD_v,\\
 H_M&=D_u(CD_v+D+\tfrac12)+AD_v+(\alpha+\tfrac12)v^2.
 \end{split}                                           \tag{RFN9}
\]
There are no omitted higher state-normal terms.

\paragraph{Exact scalar interface, before spending the $n$-mass.}
Set
\[
 U=\frac u{\sqrt x},\quad V=\frac v{\sqrt y},\quad
 Z_* =\frac z{\sqrt{q_z}},\qquad
 Q_*=U-VZ_*,\quad B_*=V-UZ_*,\quad R_*=Z_*-UV,
\]
\[
 \Psi=U^2+V^2+Z_*^2+U^2V^2Z_*^2-4UVZ_*.                 \tag{RFN10}
\]
All three normalized testing variables are positive, but none is
assumed bounded.  A legal separated certificate is
\[
 \Psi=(U-V)^2+Z_*^2(1-UV)^2+2UV(1-Z_*)^2.               \tag{RFN11}
\]
For example, $\sqrt{UV}$ is the product of the source factor
$(xy)^{-1/4}$ and the testing factor $\sqrt{uv}$.

The exact RS/RB scalar identity for the actual source is
\begin{align}
 F={}&q_cQ_*^2+nj^2y B_*^2+r_cR_*^2+d_cV^2Q_*^2
   +(np^2y^2+e_0)V^2B_*^2\\
   &+f_c\Psi+g_cV^2\Psi+k_cZ_*^2(V^2-1)^2,
 \qquad e_0=\tfrac12p^2y^2-g_c.                        \tag{RFN12}
\end{align}
For an entirely algebraic definition of the masses, put
$K_0=CH^2+(D+\tfrac12)(1-H)^2$ and
\[
\begin{array}{ll}
 M_a=Ax,&M_b=\{\mu(1-p)+(n+\tfrac12)p^2\}y,\\
 M_c=\alpha q_z,&
 M_{ab}=\{A+\alpha(1-p)^2+\mu(1-p)+\tfrac12(1-p)^2\}xy,\\
 M_{ac}=nxq_z=nj^2y,&
 M_{bc}=\{K_0+AH^2+2H\alpha+H-\tfrac12\}yq_z,\\
 M_{b^2}=(n+\tfrac12)p^2y^2,&M_{abc}=(K_0+n)xyq_z,\\
 M_{b^2c}=H^2(A+C+\alpha+\tfrac12)y^2q_z,&
 M_{ab^2c}=CH^2xy^2q_z.
\end{array}
\]
Then
\[
\begin{split}
 k_c&=(M_a-M_b+M_{ac}+M_{ab^2c}-M_{bc})/2,\\
 r_c&=M_c-M_b+M_{ac}-k_c,\quad f_c=M_b-M_{ac},\quad q_c=M_a-f_c,\\
 d_c&=M_{b^2c}-k_c,\qquad g_c=M_{ab^2c}.
\end{split}
\]
The source contact identities make (RFN12) exact, including both odd
monomials of $F$.  It is checked coefficientwise against (RFN7), not
only at the contact point.

Define the unnormalized current $b=v-uz/j=\sqrt y B_*$ and the
pure source quantities
\[
 c_b=\frac{nj}{n+h_0},\qquad M_e=\frac{nh_0j^2}{n+h_0}.
                                                               \tag{RFN13}
\]
The whole scalar pieces $nj^2b^2+np^2v^2b^2$ in (RFN12) must be kept
until their correlations with $N_0$ have been completed.  The full
$N_0$-linear term below contains $njb+npv^2b$, so these pieces and the
corresponding hard terms give
\[
 n(N_0+jb)^2+n v^2(N_0+pb)^2.
\]
Complete only the first square, using a pure source inverse:
\[
 h_0N_0^2+n(N_0+jb)^2
   =(n+h_0)(N_0+c_bb)^2+M_e b^2.                       \tag{RFN14}
\]
Keep $nv^2(N_0+pb)^2$ whole.  Put
$\widehat N=\sqrt{n+h_0}(N_0+c_bb)$.
The other hard terms in (RFN9) are the nonnegative whole squares
\[
 h_0(uN_0)^2+\tfrac12(vN_0)^2+\alpha(uN_0)^2.             \tag{RFN15}
\]
They are still kept whole after the shift; no term $nu^2$ or $nv^2$
is treated as a bounded perturbation.

For $n>0$ one has the exact bounded-shift identity
\[
 \frac{c_b^2}{M_e}=\frac{n}{h_0(n+h_0)}=O(1).
                                                               \tag{RFN16}
\]
At $n=0$, set $c_b=M_e=0$ and omit the effective current row.  Formula
(RFN14) remains literal, and no division by $n$ or $\sqrt{M_e}$ is
made in that case.

RS pays the signed $e_0$ and $k_c$ terms in (RFN12) at infinitesimal
cost after (RFN14), with the effective current mass $M_e y$, and leaves
fixed positive real fractions of
\[
 q_cQ_*^2,\qquad M_e b^2,\qquad r_cR_*^2,\qquad f_c\Psi. \tag{RFN17}
\]
The second term is omitted when $n=0$.  This is the precise scalar
allocation needed here; an allocation that spends the raw $nj^2y$
mass before (RFN14) would not suffice.

\subsection{Finite exact budgets for the full normal rows}

In the rest of the proof, a row is called free if any prescribed fixed
real multiple of its square can be purchased at infinitesimal cost
from (RFN17).  This has an exact meaning: a finite SOS identity writes
its square below a finite sum of those retained squares, with bounded
separated coefficients in $\mathfrak n$.  Subtracting any fixed real
multiple changes their strict Gram diagonals only infinitesimally.

\paragraph{Three elementary circuit budgets.}
Direct expansion gives
\begin{align}
 (1+U^2)\Psi-R_*^2
     &=B_*^2+U^2Q_*^2+U^2(1-UVZ_*)^2,                  \tag{RFN18a}\\
 (1+V^2)\Psi-V^2B_*^2
     &=Q_*^2+R_*^2+V^2(1-UVZ_*)^2.                   \tag{RFN18b}
\end{align}
There is also a useful opposite-current budget.  Put
$t_*=UV$, $\delta_*=U-V$, $E_*=Z_*(1-t_*)$,
$T_*=t_*(1-Z_*)$, and $J_*=T_*-\delta_*^2Z_*$.  Then
\begin{align}
 2(\Psi+V^2Q_*^2+U^2B_*^2)-R_*^2
 ={}&(E_*+T_*)^2+2\delta_*^2+4t_*(1-Z_*)^2\\
   &+2J_*^2+4t_*\delta_*^2Z_*^2.                      \tag{RFN18c}
\end{align}
Every term on the right is a legal separated square.  Thus a
nonnegative multiplier $a$ buys $aR_*^2$ at infinitesimal cost whenever
\[
 a/f_c\in\mathfrak n,\qquad a/(q_cy)=O(1),\qquad
 a/(M_e yx)=O(1),                                      \tag{RFN18d}
\]
since the last two diagonal costs have the additional testing factors
$v^2$ and $u^2$.  This criterion is used only when $M_e>0$.

\paragraph{An exact radial budget.}
The row $R_u=v(u^2-x)$ is free in all source subcases, including the
weak opening $1-p\in\mathfrak m_c$.  Here is an explicit finite
certificate, so that this assertion does not rest on pointwise
nonnegativity.  Write $D_*=V(U^2-1)$ and
\[
\begin{split}
 t_7={}&(Z_*-U^2V,\ V-U^2V,\ U-U^2V,\ UZ_*-U^2V,\\
       &\hspace{25mm}UV-U^2V,\ UVZ_*-U^2V,\ U^2-U^2V)^T,\\
 t_3={}&(1-UV,\ V-UV,\ U-UV)^T,
\end{split}
\]
\[
 G_7=\begin{pmatrix}
 4&-2&-3/2&0&-1/2&0&1/2\\
 -2&3&0&-3/2&0&-1&2\\
 -3/2&0&4&3/2&-2&-3&0\\
 0&-3/2&3/2&4&-3/2&0&-3/2\\
 -1/2&0&-2&-3/2&3&2&-1\\
 0&-1&-3&0&2&8&-4\\
 1/2&2&0&-3/2&-1&-4&4
 \end{pmatrix},\qquad
 G_3=\begin{pmatrix}3&-2&-2\\-2&2&2\\-2&2&3\end{pmatrix}.
\]
The exact identity is
\[
 4\Psi+4U^2Q_*^2+4U^2R_*^2-D_*^2
   =t_7^TG_7t_7+4VZ_*(1-U)^2+UZ_*t_3^TG_3t_3.         \tag{RFN19}
\]
The leading principal minors of $G_7$ are
$4,8,101/4,1049/16,907/16,3279/16,743/16$; those of $G_3$ are
$3,2,2$.  Both are positive definite real matrices.  Their real
Cholesky factors and the separated square roots of $VZ_*$ and $UZ_*$
therefore make the right side a finite legal SOS.  Denote it by
$\mathcal W(U,V,Z_*)$.

Since $R_u=x\sqrt y D_*$, the exhaustive source cases are as follows.

\emph{Case 1: $1-p$ is a positive unit.}
Multiply (RFN19) by $x^2y$.  The costs relative to
$f_c\Psi,q_cQ_*^2,r_cR_*^2$ are, respectively,
\[
 O(B),\qquad O(Bu^2),\qquad O(u^2/(1-p)).
\]
All belong to $\mathfrak n$.

\emph{Case 2: $p\sim1$ and $\eta_*$ is a positive unit.}
Then $j$ has a positive real lower bound, and $j^2B$ is a positive
unit.  Adding $4U^2$ times (RFN18a) to (RFN19) gives
\[
 D_*^2\ \preceq_{\rm sos}\
 4(1+U^2+U^4)\Psi+4U^2Q_*^2.
\]
After multiplication by $x^2y$, the costs are
\[
 O\!\left(B+u^2/j+u^4/(j^2B)\right),\qquad O(Bu^2),
\]
again infinitesimal.  The notation $a\preceq_{\rm sos}b$ means
$b-a\in\Sigma\mathcal T^2$ and is only shorthand for the displayed
exact identities.

\emph{Case 3: $p\sim1$ and $\eta_*\in\mathfrak m_c$.}
Formula (RFN3d) gives $n>0$ and $1/M_e=O(1)$; it also gives
$p^2/M_e=O(1)$.  Put $E^{\rm rad}=Z_*(U^2-1)$.
Permuting $V$ and $Z_*$ in (RFN19) yields
\[
 \mathcal W(U,Z_*,V)
   =4\Psi+4U^2Q_*^2+4U^2B_*^2-(E^{\rm rad})^2.
\]
The identity $D_*=UE^{\rm rad}+(U^2-1)B_*$ then gives, exactly,
\begin{align}
 &8U^2\Psi+8U^4Q_*^2+
   \{8U^4+2(U^2-1)^2\}B_*^2-D_*^2\\
 &\hspace{10mm}=2U^2\mathcal W(U,Z_*,V)
                 +\{UE^{\rm rad}-(U^2-1)B_*\}^2.        \tag{RFN20}
\end{align}
After multiplication by $x^2y$, the three costs are
\[
 O(u^2/j),\qquad O(u^4/j),\qquad
 \frac{8u^4+2(u^2-x)^2}{M_e}.
\]
They are infinitesimal.  These three cases prove the asserted free
radial row without assuming a rate for any testing variable.

\paragraph{The smaller current helpers.}
If $\eta_*\in\mathfrak m_c$, (RFN3d) implies
\[
 \frac{p^2}{M_e}
 =\frac{p^2}{h_0j^2}+\frac{p^2}{nj^2}=O(1),
\]
because $p/j<1$ and
$p(j-p)(j+1)/j^2\leq p+p/j$.
Thus $(pv^2b)^2$ costs only $(p^2/M_e)v^4$ of the effective current.
If $\eta_*$ is a positive unit instead, then
\[
 \frac{p^2y}{f_c}=O\!\left(\frac{p+p/j}{\eta_*}\right)=O(1).
\]
Multiply (RFN18b) by $p^2y^2v^2$.  This buys $(pv^2b)^2$ with circuit
cost
\[
 \frac{p^2y}{f_c}(yv^2+v^4)\in\mathfrak n.
\]
The same identity multiplied by $Tp^2u^2y^2$ shows that
$\sqrt T\,puvb$ is free: its cost is either
$(p^2/M_e)Tu^2v^2$ or
$(p^2y/f_c)Tu^2(y+v^2)$.

\subsection{The entire two-normal conductor}

\paragraph{First normal: exact source cancellation and allocation.}
Put
\[
 q=vz-Tu,\qquad R=z-uv/B,\qquad R_u=v(u^2-x),
\]
\[
 K_c=\frac p2+j\{H(\alpha+\tfrac12)-\tfrac12\}+n(p-j),
 \qquad
 K_R=-\mu-(\alpha+\tfrac12)(1-p)+\frac TjK_c.
\]
Both $K_c/j$ and $K_R$ are bounded by (RFN3).  The complete first normal
is exactly
\[
 L_N=\{j(n-\alpha)+p(n+\tfrac12)v^2\}b
                  +K_RR_u+\frac{K_c}{j}uvq.            \tag{RFN21}
\]
In particular the sign of the $\mu$ term in $K_R$ is negative.
For a direct source check, the difference between this formula and
the derivative of (RFN7) is $vS_N$, where
\[
 S_N=np-\mu+\tfrac p2+
       \{-\mu-(\alpha+\tfrac12)(1-p)\}x
       +j(\alpha-n)+yK_c=0.
\]
The last equality is an exact rational identity of the actual native
source under (RFN2), not an asymptotic contact equation.

After the two correlated $n$-squares are formed, the remaining linear
coefficient is
\[
 L_N^{\rm rem}=-j\alpha b+\tfrac12pv^2b
                       +K_RR_u+\frac{K_c}{j}uvq.
\]
The $q$ row is strict because
$q_cQ_*^2=(q_c/(Tx))(q^2/T)$, with $q_c/(Tx)$ a positive unit.
Its coefficient in $L_N^{\rm rem}$ relative to $q/\sqrt T$ is
$(K_c/j)uv\sqrt T\in\mathfrak n$.
The $pv^2b$ helper is free by the preceding paragraph.
For the remaining small row, the exact identity
\[
 j\alpha b=-\frac\alpha B R_u-\alpha uR
\]
has $\alpha/B=O(B+x)\in\mathfrak m_c$.
When $1-p$ is a unit, the square of $\alpha uR$ costs
$O(B^2(j+1)^2u^2/(1-p))\in\mathfrak n$ against $r_cR_*^2$.
When $p\sim1$, use $B_*^2\preceq_{\rm sos}(1+U^2)\Psi$ from
(RFN18a).  The circuit cost of $(j\alpha b)^2$ is
\[
 O\!\left(B^3(j+1)^2+B^2(j+1)^2u^2/j\right)\in\mathfrak n,
\]
since $j$ is bounded below and $B(j+1)=B+x$ is infinitesimal.
Finally, $K_RR_u$ has bounded coefficient; buy an arbitrarily large
fixed real $R_u^2$ pivot using (RFN19)--(RFN20).

The shift $N_0=\widehat N/\sqrt{n+h_0}-c_bb$ introduces no new
unbounded Gram coefficient: $1/\sqrt{n+h_0}=O(1)$, and for $n>0$
$c_b/\sqrt{M_e}=O(1)$ by (RFN16).  The only bounded, potentially
non-infinitesimal couplings are therefore those of $R_u$ to
$\widehat N$ and $\sqrt{M_e}b$.  A sufficiently large fixed real radial
pivot absorbs both while preserving a strict positive real residue.
For $n=0$ the latter coupling is exactly absent.

\paragraph{Second normal: cancellation of both source openings.}
Define
\[
 K_d=CH-(D+\tfrac12)(1-H),\qquad
 C_q=K_d+AH+\tfrac12
       +v^2\{CH(1+u^2)+AH+(\alpha+\tfrac12)H\}.
\]
The complete second normal is
\[
 L_M=C_q q+\mu uvb+v(\alpha-BTCHu^2)R.                  \tag{RFN22}
\]
The actual source satisfies the three exact identities
\[
\begin{split}
 S_0&=T(K_d+AH+\tfrac12)-A=0,\\
 S_1&=jK_d+\mu+xTCH=0,\\
 S_2&=jK_d+xTCH+T\{CH+AH+(\alpha+\tfrac12)H\}
          -A-(\alpha+\tfrac12)(1-p)+\alpha/B=0.
\end{split}
\]
Before imposing these identities, the complete derivative minus
(RFN22) is exactly
$uS_0+u^2vzS_1/j+uv^2(S_2-S_1)$.
This records all state terms and rules out a hidden $j$-weighted
radial error.  In particular $S_1$ gives
$K_d=-\mu/j-BTCH=O(1)$ even when $C$ or $D$ is unbounded.
Consequently $TC_q=O(T+v^2)\in\mathfrak n$.

Normalize
\[
 M_0=\sqrt{TH(1-H)}\,\widehat M.
\]
The full diagonal (RFN9), not only its zero-angle value, obeys
\[
 TH(1-H)H_M\sim k/2>0.                                 \tag{RFN23}
\]
Indeed the $C$ and $D$ terms contribute
$(k/2)D_u\{(1-H)D_v+H\}$ up to infinitesimal bounded coefficients,
and all remaining terms are infinitesimal after this normalization.
This is uniform at both independent openings.
The normalized $C_qq$ coupling has coefficient
$T\sqrt{H(1-H)}C_q\in\mathfrak n$ on the strict row $q/\sqrt T$.
The normalized $\mu uvb$ row is free because $\mu=O(p)$ and
$\sqrt T\,puvb$ has already been purchased.

It remains to allocate the entire opposite row, including its
higher testing terms.  Put
\[
 J=B(j+1)+u^2\in\mathfrak n,\qquad
 a_R=Tq_zv^2(\alpha-BTCHu^2)^2.
\]
The factor $H(1-H)$ may be discarded from the cost bound, since it is
bounded.  Since $\alpha-BTCHu^2=O(BJ)$, the exact squared multiplier
$a_R$ satisfies
\[
 \frac{a_R}{f_c}=O((T/j)v^2J^2),\qquad
 \frac{a_R}{q_cy}=O(v^2J^2).                            \tag{RFN24a}
\]
If $1-p$ is a unit, its direct cost against $r_cR_*^2$ is
$O(Tv^2J^2/(1-p))\in\mathfrak n$.
If $p\sim1$ and $\eta_*$ is a unit, apply (RFN18a); the cost is
$(a_R/f_c)(1+u^2/x)\in\mathfrak n$, using that $j^2B$ is a unit.
If $p\sim1$ and $\eta_*\in\mathfrak m_c$, (RFN3d) gives
\[
 \frac{a_R}{M_e yx}
     =O\!\left(\frac{T(1+1/n)}{j^2}v^2J^2\right)=O(1),  \tag{RFN24b}
\]
because $(1+1/n)/j^2=O(1)$ in this source case.  The other two
conditions in (RFN18d) follow from (RFN24a).  The exact budget
(RFN18c) therefore buys this row at infinitesimal cost as well.
These estimates treat the whole square
$(\alpha-BTCHu^2)^2$; there is no truncation at a testing degree.

\paragraph{Full Gram, legal roots, and restoration.}
After the normalizations, the entire remaining mixed hard term is
\[
 2(\alpha+\tfrac12)uv\sqrt{TH(1-H)}
    \left(\frac{\widehat N}{\sqrt{n+h_0}}-c_bb\right)\widehat M.
\]
Its coefficients relative to $\widehat N$, $\widehat M$, and, when
present, $\sqrt{M_e}b$ are infinitesimal; (RFN16) controls the shift.
All other remaining normal couplings were allocated above.  Choose a
fixed sufficiently large radial pivot, and then purchase all the
finitely many free helpers.  Their scalar costs lie in $\mathfrak n$,
so fixed positive fractions of (RFN17) survive.  The complete bounded
Gram has strict positive real residue: the $\widehat N$ pivot is one,
the $\widehat M$ pivot has residue $\overline{k}/2$, the retained scalar pivots
are positive, and the only bounded radial off-diagonals have been
absorbed by the chosen real radial pivot.

The separated strict bounded-Gram lemma now gives a finite sum of
squares in $\mathcal T$.  Its row entries need not be bounded.  All
budget remainders are explicit finite separated sums of squares, all
source-only completions use positive pure source denominators, and
all state-coordinate inverses are pure testing.  Restore the whole
squares (RFN15), the retained $nv^2(N_0+pb)^2$ square, the scalar
budget remainders, the $mpD_u^2N_0^2$ packet, and the pure testing
factor $D_u$.  This certifies the complete original parent (RFN4).
Every bounded coefficient used in the argument depends only on the
source and on $u,v$.  The arbitrary state ratio $z$ remains inside
retained rows throughout.  Hence the proof has no bounded-state
restriction, and it includes literal $n=0$ without any limiting
argument.

\paragraph{Audit boundary.}
The exact checker validates the complete parent, both full normals,
all hard coefficients, the actual source cancellations, the scalar
circuit, and the finite Gram/current budget identities.  The source
bounds and the signed scalar allocation are the separately stated RS
input; the analytical argument above proves the full normal lift
from that input.  A separate state/sign proof and a separate outer
angular join remain necessary for a whole-source all-testing claim.
No explicit global LMI, optimized lift size, or external formal proof
is claimed here.
